% =============================================================================
% Marco Conceptual
% =============================================================================
\section{Marco Conceptual}

\subsection{Conceptos Fundamentales}

\subsubsection{Redes neuronales artificiales (ANN)}

 Modelos computacionales inspirados en el funcionamiento del cerebro humano, su elemento principal son las neuronas. Las ANN se caracterizan por una asamblea de cientos de estas neuronas artificiales (unidades de computación básicas). El entrenamiento consiste en un par de señales como entrada y también se les presentan una serie de salidas deseadas, imitando como el cerebro depende de estímulos externos para aprender.
 
\subsubsection{Redes Neuronales Convolucionales (CNN)}
 
Las CNN están diseñadas para imitar la corteza visual humana y procesar datos bidimensionales o tridimensionales. Las CNN son la evolución de las ANN para procesamiento de imágenes. Su importancia técnica radica en:

\begin{itemize}
    \item Eliminación de la necesidad de extraer características manualmente
    \item Aprendizaje de jerarquías de patrones (bordes, líneas, formas, texturas y patologías)
    \item Procesamiento eficiente de datos espaciales
\end{itemize} 

\paragraph{Arquitectura y Capas Principales}

El funcionamiento de las CNN está compuesto por varias capas especializadas \parencite[p.~28]{Seeram2024}:

\begin{enumerate}
    \item \textbf{Capa de entrada}: Recibe la imagen, donde los píxeles se consideran neuronas.
    \item \textbf{Capa de convolución}: Filtros que se deslizan sobre la imagen generando mapas de características.
    \item \textbf{Función de Activación (ReLU)}: Introduce no linealidad convirtiendo valores negativos en cero.
    \item \textbf{Capa de agrupación (Pooling)}: Reduce la dimensionalidad conservando las características más importantes.
    \item \textbf{Capas totalmente conectadas}: Unifican la información para razonamiento de alto nivel.
    \item \textbf{Capa de salida}: Asigna una probabilidad a cada categoría diagnóstica mediante funciones como \enquote{Softmax}.
\end{enumerate}

Durante el entrenamiento, la imagen pasa por todas las capas generando una predicción que se compara con el resultado esperado, proporcionado por radiólogos expertos; los errores se propagan de nuevo en la red para ir ajustando la salida hasta que el error sea mínimo.

\subsection{ImageNet}
La base de datos más grande del mundo para el entrenamiento de redes neuronales convolucionales, contiene más de 1.2 millones de imágenes naturales clasificadas en más de 1000 categorías. ImageNet es la base para una competición anual ImageNet Large Scale Visual Recognition Challenge (ILSVRC), la cual se ha convertido en el estándar para evaluar el rendimiento de los algoritmos de reconocimiento de objetos a gran escala.

Los investigadores utilizan ImageNet para una técnica llamada transferencia de aprendizaje (transfer learning), que consiste en utilizar un modelo ya entrenado con el millón de imágenes de ImageNet. Como este modelo ya ha aprendido a detectar bastantes características fundamentales como bordes, texturas y formas, este modelo ya inteligente se adapta a una tarea médica específica (detectar tumores, neumonía, etc.), todo esto ahorra tiempo y recursos computacionales.


\subsubsection{Detección Asistida por Computadora (CADe)}

Su objetivo principal es identificar y localizar anomalías en imágenes médicas. Se enfoca especialmente en reducir las tasas de falsos negativos, es decir, identificar tumores o lesiones que podrían pasar desapercibidos para el ojo humano.
  
\subsubsection{Diagnóstico Asistido por Computadora (CADx)}
Esto va un paso más allá de la simple detección, enfocándose en la evaluación y caracterización de las patologías ya identificadas. Por ejemplo, analiza la salida radiográfica para evaluar la probabilidad de que una lesión sea benigna o maligna. También intenta especificar el tipo de enfermedad, su severidad, estadio o la intervención recomendada.
 
\subsection{Arquitecturas Relevantes}

\begin{description}
    \item[AlexNet] La red que inició el interés masivo al superar a todos los algoritmos en la competencia ImageNet. Se considera una arquitectura poco profunda comparada con las actuales, pero fue la que demostró que entre más capas, mayor precisión (hasta cierto punto). Se utiliza mayormente para clasificación de imágenes, donde se asigna una etiqueta a la imagen (¿hay tumor o no?).
    
    \item[U-Net] La base perfecta para la segmentación de imágenes médicas debido a su estructura simétrica en forma de U y \enquote{conexiones de salto} que recuperan información espacial precisa. A diferencia de AlexNet, que predice una etiqueta para toda la imagen, U-Net realiza una predicción píxel por píxel, lo que permite delinear el contorno exacto de un objeto. Es la preferida para segmentación (proceso de identificar y delimitar con precisión píxeles) de imágenes. 
    
    \item[ResNet] AlexNet probó que entre más capas (más profundidad) mayor sería la precisión. Sin embargo, con modelos como VGGNet o GoogLeNet se observó que al aumentar la profundidad surgieron dos problemas: el desvanecimiento del gradiente y la degradación de la precisión. ResNet introdujo conexiones residuales que permiten que la información \enquote{salte} capas. Esto permitió entrenar redes extremadamente profundas sin pérdida de precisión. Se utiliza en tareas de alta precisión y reconocimiento de patrones complejos (reconocimiento de melanoma, radiografías de tórax, etc.).
\end{description}

\subsection{Métricas}
La precisión por sí sola puede ser engañosa en diagnósticos médicos. Un modelo que siempre marca a los pacientes como sanos tendría alta precisión si la mayoría lo están, pero fallaría en detectar a los enfermos. Por ello se utilizan métricas complementarias basadas en la matriz de confusión: verdaderos positivos (VP), verdaderos negativos (VN), falsos positivos (FP) y falsos negativos (FN).

\paragraph{Sensibilidad} Mide qué tan bien el modelo identifica los casos positivos:
\[
\text{Sensibilidad} = \frac{\text{VP}}{\text{VP} + \text{FN}}
\]

\paragraph{Especificidad} Mide la capacidad de etiquetar correctamente a pacientes sanos:
\[
\text{Especificidad} = \frac{\text{VN}}{\text{VN} + \text{FP}}
\]

Ambas métricas se integran en la Curva ROC (Receiver Operating Characteristic), cuya área bajo la curva (AUC) resume la eficiencia del modelo en un valor único entre 0 y 1. Un AUC de 1.0 representa la perfección.